\begin{abstract}
Recent advancements exemplified by BitNet~\cite{bitnet} are ushering in a new era characterized by 1-bit Large Language Models (LLMs). In this study, we propose a ternary 1-bit LLM, \textbf{\text{BitNet b1.58}}, where each parameter or weight of the LLM assumes a value from the set \{-1, 0, 1\}. This model attains comparable perplexity and downstream task performance to its full-precision (FP16 or BF16) Transformer counterparts with identical model sizes and number of training tokens, while offering marked improvements in efficiency regarding memory usage, latency, throughput, and energy consumption. Importantly, the 1.58-bit LLM establishes a revised scaling law and methodology for developing future generations of LLMs that reconcile strong performance with enhanced efficiency. In addition, this approach facilitates new computational paradigms and motivates the design of specialized hardware that takes advantage of the unique properties of 1-bit LLMs.
\end{abstract}

\section{The Era of 1-bit LLMs}

The AI landscape has recently experienced explosive advancements in the scale and abilities of Large Language Models (LLMs). Although these systems achieve outstanding results across diverse natural language processing tasks, the ever-increasing number of model parameters introduces deployment bottlenecks and raises issues of sustainability and cost efficiency due to substantial energy demands.

A prominent strategy to mitigate these drawbacks involves the use of post-training quantization techniques, which compress models into lower-precision representations for inference~\cite{smoothquant, gptq, quip, quip_sharp}. By quantizing both weights and activations, these methods considerably decrease the required computational and memory resources for LLM operation. While this practice has driven reductions in bit-widths from 16-bit to 4-bit formats~\cite{gptq, awq}, and is common in industrial deployments, post-training quantization still suffers from inherent sub-optimality.

Recent advances in 1-bit model designs, exemplified by BitNet~\cite{bitnet}, provide a compelling approach for minimizing the computational costs associated with LLMs without sacrificing effectiveness. Standard LLMs typically utilize 16-bit floating-point representations (such as FP16 or BF16), with matrix multiplication constituting the core computational workload. Consequently, the principal source of computational expense arises from operations involving floating-point multiplication and addition. By contrast, BitNet employs matrix multiplication relying solely on integer addition, offering significant energy savings when compared to traditional methods. Given that the limiting factor for computational throughput in many hardware systems is power consumption, these efficiency gains also yield potential improvements in computation speed.

Beyond the computational overhead, moving model weights from off-chip DRAM to on-chip memory (such as SRAM) during inference is another costly component. Efforts to increase on-chip SRAM to alleviate this bottleneck often result in far higher costs relative to DRAM. With a far smaller memory requirement, both in terms of capacity and bandwidth, 1-bit LLMs significantly reduce both the expense and latency associated with DRAM-to-accelerator weight transfers. This leads to substantial acceleration and resource savings in the inference process.

In this study, we present a novel 1-bit LLM variant termed \textbf{\text{BitNet b1.58}}, where each parameter is chosen from the ternary set \{-1, 0, 1\}. By introducing 0 to the initial 1-bit BitNet parameterization, the effective representation increases to 1.58 bits per parameter in binary terms. \text{BitNet b1.58} maintains all of the advantages observed in the original BitNet, such as an efficient computational model that virtually eliminates multiplication in matrix operations and is highly amenable to optimization. Its energy efficiency remains on par with the initial 1-bit architecture and is markedly superior to FP16 LLMs in terms of memory use, throughput, and latency. Moreover, \text{BitNet b1.58} confers two further benefits: first, it improves representational power through explicit support for feature selection—made feasible by the inclusion of 0-valued weights—which can notably enhance 1-bit LLM capabilities; second, experimental results demonstrate that \text{BitNet b1.58}, beginning at a 3B parameter scale and under matched configuration settings (model size, token count, etc.), achieves parity with full-precision (FP16) baselines in perplexity and downstream task accuracy.

\section{BitNet b1.58}

\text{BitNet b1.58} extends the BitNet architecture, which is a Transformer variant substituting \emph{nn.Linear} layers with \emph{BitLinear}. The model is trained from initialization using weights quantized to 1.58 bits and activations quantized to 8 bits. While rooted in the original BitNet, it incorporates several noteworthy changes, summarized below.

\paragraph{Quantization Function.} To limit the weights to the discrete values of -1, 0, or +1, an \emph{absmean} quantization scheme is employed. This technique normalizes the weight matrix by dividing it by its mean absolute value before quantizing each element to the closest integer in \{-1, 0, +1\}:

\begin{equation}
    \widetilde{W} = \mathrm{RoundClip}(\frac{W}{\gamma + \epsilon}, -1, 1),
\end{equation}

\begin{equation}
    \mathrm{RoundClip}(x, a, b) = \max(a, \min(b, \mathrm{round}(x))),
\end{equation}

\begin{equation}
    \gamma = \frac{1}{nm}\sum_{ij} |W_{ij}|.
\end{equation}

The process for activation quantization adheres to the method implemented in BitNet, with one key distinction: instead of normalizing activations before non-linearities to $[0, Q_b]$, all activations are scaled to the interval $[-Q_b, Q_b]$ on a per-token basis, thereby omitting zero-point quantization. This approach streamlines both the implementation and system-level optimization, and our experimental results indicate that the impact on model performance is minimal.

\paragraph{LLaMA-alike Components.}

LLaMA~\cite{llama, llama2} has become the reference architecture for open-source large language models. To facilitate compatibility with this ecosystem, \text{BitNet b1.58} incorporates key LLaMA-inspired modules. Notably, it employs RMSNorm~\cite{rmsnorm}, SwiGLU~\cite{swiglu}, rotary embeddings~\cite{rope}, and omits all bias parameters. These architectural choices ensure that \text{BitNet b1.58} can be seamlessly integrated with widely-used open-source libraries and frameworks, including Huggingface, vLLM~\cite{vllm}, and llama.cpp\footnote{\url{https://github.com/ggerganov/llama.cpp}}.

\section{Results}

We conducted a comparison between \text{BitNet b1.58} and our replicated FP16 \text{LLaMA LLM} across different parameter scales. Both sets of models were pre-trained on the RedPajama dataset~\cite{redpajama}, each for 100 billion tokens to maintain experimental parity. For evaluation, we utilized a variety of zero-shot language understanding benchmarks, including ARC-Easy~\cite{arc}, ARC-Challenge~\cite{arc}, Hellaswag~\cite{hellaswag}, Winogrande~\cite{winoGrande}, PIQA~\cite{piqa}, OpenbookQA~\cite{openbookqa}, and BoolQ~\cite{boolq}. Additionally, validation perplexity metrics were assessed on the WikiText2~\cite{wikitext2} and C4~\cite{c4} corpora.

To assess resource requirements and efficiency, we measured GPU memory consumption and inference latency for both \text{LLaMA LLM} and \text{BitNet b1.58}. All timing and memory measurements were collected using the FasterTransformer\footnote{\url{https://github.com/NVIDIA/FasterTransformer}} library, which provides highly optimized LLM inference on GPUs. For \text{BitNet b1.58}, we also incorporated the 2-bit Ladder kernel~\cite{ladder}. Our primary inference efficiency metric was the elapsed time per generated output token, which dominates inference cost.

A summary of perplexity scores and associated resource costs for both \text{BitNet b1.58} and \text{LLaMA LLM} is provided in Table~\ref{tab:ppl}. The findings indicate that at the 3B parameter scale, \text{BitNet b1.58} achieves perplexity comparable to the full-precision \text{LLaMA LLM}, while requiring only 37\% of the GPU memory and running 2.71 times faster. Notably, the 3.9B configuration of \text{BitNet b1.58} delivers even better results, being 2.4 times faster and utilizing just 30\% of the GPU memory of the \text{LLaMA LLM} 3B, yet substantially outperforming it.

Detailed zero-shot task performance is presented in Table~\ref{tab:zero-shot}. Evaluation protocols followed those of \emph{lm-evaluation-harness}\footnote{\url{https://github.com/EleutherAI/lm-evaluation-harness}}. The data reveal that as model scale increases, the difference in accuracy between \text{BitNet b1.58} and \text{LLaMA LLM} decreases. Crucially, at and above the 3B scale, \text{BitNet b1.58} achieves parity with the full-precision baseline. Echoing the perplexity trends, end-task accuracy results demonstrate that \text{BitNet b1.58} 3.9B exceeds the performance of \text{LLaMA LLM} 3B, while demanding less memory and offering faster inference. Collectively, these results establish \text{BitNet b1.58} as a Pareto improvement over leading large language models.

\paragraph{Memory and Latency}

We extended our experiments to include larger model configurations—7B, 13B, and 70B parameters—and analyzed the associated computational costs. The patterns depicted in Figure~\ref{fig:latency-memory-energy} indicate that latency and memory requirements both trend upward with model scaling, while the overall acceleration improves for bigger models. Notably, \text{BitNet b1.58} 70B demonstrates a speedup factor of 4.1 compared to the \text{LLaMA LLM} reference. This improvement stems from the increasing contribution of \emph{nn.Linear} to runtime as the parameter count grows. Memory usage exhibits a parallel trajectory; with embeddings maintained at full precision, their relative impact on memory declines in larger architectures. All benchmarks utilized a 2-bit kernel, implying potential headroom for future cost reductions through further optimizations.

\paragraph{Energy}

We also performed an assessment of the arithmetic energy demands for both \text{BitNet b1.58} and \text{LLaMA LLM}. Matrix multiplication operations were our principal focus, since these account for the predominant portion of LLM computation costs. As shown in Figure~\ref{fig:energy}, \text{BitNet b1.58} expends most of its energy on INT8 addition operations, while \text{LLaMA LLM} mainly involves FP16 addition and multiplication. Relying on the energy estimation frameworks from~\cite{energycost, pokebnn}, our findings indicate that \text{BitNet b1.58} reduces energy expenditure for matrix multiplications by a factor of 71.4 when operating on 7nm fabrication processes. Additionally, we computed the overall energy required to process sequences of 512 tokens end-to-end. The outcomes demonstrate a growing energy efficiency advantage for \text{BitNet b1.58} over the FP16 \text{LLaMA LLM} baseline as the parameter count increases. This is attributed to the fact that the proportional overhead from \emph{nn.Linear} intensifies with scale, while contributions from other components diminish in significance in larger networks.

\paragraph{Throughput}

We evaluate and contrast the throughput capabilities of \text{BitNet b1.58} and \text{LLaMA LLM} with 70B parameters, utilizing two 80GB A100 GPUs and implementing pipeline parallelism~\cite{gpipe} to ensure the feasibility of running the \text{LLaMA LLM} 70B on this hardware. Batch size was increased incrementally until the GPU memory capacity was reached, with a sequence length set at 512. According to the results summarized in Table~\ref{tab:throughput}, \text{BitNet b1.58} 70B supports batch sizes up to 11 times greater than those accommodated by \text{LLaMA LLM}, culminating in an 8.9-fold improvement in throughput.

\textbf{\text{BitNet b1.58} introduces a novel scaling law relating model efficacy to inference cost}. By referencing the data presented in Figure~\ref{fig:latency-memory-energy} and \ref{fig:energy}, we can establish the following equivalencies between model sizes at 1.58 bits and their 16-bit FP16 counterparts:

\begin{itemize}
    \item 13B BitNet b1.58 surpasses 3B FP16 LLM in terms of latency, memory consumption, and energy efficiency.
    \item 30B BitNet b1.58 demonstrates greater efficiency across latency, memory usage, and energy consumption than 7B FP16 LLM.
    \item 70B BitNet b1.58 is more efficient with respect to latency, memory, and energy requirements than 13B FP16 LLM.
\end{itemize}

\paragraph{Training with 2T Tokens}

The total number of tokens used for training remains a pivotal consideration for LLM scalability. To examine the token scalability of \text{BitNet b1.58}, we conducted training using 2T tokens, adhering to the StableLM-3B~\cite{StableLM-3B-4E1T} data methodology, which reflects state-of-the-art practice for open-source 3B models. We assessed both models on a composite benchmark comprising Winogrande~\cite{winoGrande}, PIQA~\cite{piqa}, SciQ~\cite{sciq}, LAMBADA~\cite{lambada}, and ARC-easy~\cite{arc}. Zero-shot accuracy was recorded as indicated in Table~\ref{tab:2t}. For benchmarks using both accuracy and normalized accuracy, we computed the mean value. The StableLM 3B results at 2T tokens are reproduced directly from their technical report. Our analysis reveals that \text{BitNet b1.58} consistently delivers better performance across all evaluated tasks, demonstrating that 1.58-bit LLMs possess robust generalization performance.

\section{Discussion and Future Work}

\textbf{1-bit Mixture-of-Experts (MoE) LLMs}

The Mixture-of-Experts (MoE) paradigm has emerged as an efficient method for scaling LLMs in a cost-conscious manner. While MoE architectures considerably cut the required computational FLOPs, their significant demands on memory and inter-device communication bandwidth create barriers to practical adoption. These obstacles can be alleviated through the implementation of 1.58-bit LLMs. By shrinking the memory usage, fewer computational devices become necessary for MoE deployment. Additionally, this compression markedly diminishes the cost of shuttling activation data between nodes in distributed systems. Ideally, such savings could eliminate communication costs entirely, should the model fit onto a single processor.

\textbf{Native Support of Long Sequence in LLMs}

With the rise of LLMs, the need for native long sequence processing has intensified. A primary impediment in long sequence inference arises from the substantial memory needs dictated by KV caches. \text{BitNet b1.58} presents a notable advancement by cutting the precision of activation storage from 16 bits to 8 bits, thereby permitting twice the sequence length within fixed hardware constraints. Furthermore, future directions include the potential to compress these further—possibly down to 4 bits or less—without introducing any loss for 1.58-bit LLMs, an area we identify for continued exploration.

\textbf{LLMs on Edge and Mobile}

Deploying 1.58-bit LLMs offers transformative possibilities for running language models directly on edge and mobile hardware. These platforms are typically constrained by memory and compute, factors which limit the operational efficiency and size of resident LLMs. Leveraging the lower memory and energy profile of 1.58-bit LLMs opens up previously unattainable applications on such devices, considerably broadening their functional potential. Beyond enhanced feasibility, these models align well with CPU-based architectures, which predominate in edge and mobile environments. As a result, \text{BitNet b1.58} can deliver both higher efficiency and superior capability in these contexts.

\textbf{New Hardware for 1-bit LLMs}

Innovations such as Groq\footnote{\url{https://groq.com/}} have illustrated the promise of specialized hardware solutions (for instance, LPUs) tailored to LLMs. Looking ahead, we advocate for and anticipate the design and development of novel hardware architectures and system software specifically optimized for 1-bit LLMs, capitalizing on the unique computation model introduced by BitNet~\cite{bitnet}.